\documentclass[twocolumn]{article}
\usepackage{amsmath}
\title{A Tiny Example of Partial Credit}
\author{Example Author}
\date{}
\begin{document}
\maketitle
\section{Scoring rule}\label{sec:rule}
Score only items with observed evidence. The partial credit for a nonempty
set of scored items is their mean:
\begin{equation}
  s = \frac{1}{n}\sum_{i=1}^{n} x_i .
\end{equation}
Table~\ref{tab:example} shows two observed items.
\begin{table}[h]
\centering
\caption{Scores from two observed items.}\label{tab:example}
\begin{tabular}{lr}
Item & Score\\\hline
Setup & 1\\
Validation & 0\\
\end{tabular}
\end{table}
\section{Interpretation}
Items without evidence are omitted from this example's calculation.
By Section~\ref{sec:rule}, the two observed items give $s=0.5$.
\end{document}
